\section{Authenticated Encryption Security}
\label{sec:aead-security}

This section proves the AEAD security of \TW{}. The first theorem is a
random oracle multi-user indistinguishability bound. The public-permutation
theorem then follows by the flat sponge lift of \Cref{app:flat-sponge-lift},
and the all-in-one AE bound is the single-user corollary. The section then
proves \textsf{INT-RUP} integrity in the separated-interface model of
\cite{ABLMMY14}. The root tag hash and committing results are proved separately in
\Cref{sec:root-tag-hash-security}; all concrete instantiations are summarized in
\Cref{sec:concrete-security}.

All random oracle adversaries in this section are flat-RO adversaries
(\Cref{sec:ro-model}) with the displayed resources, and the flat sponge lift
transfers each random oracle bound using the derived-adversary cap
$\qRO(\mathcal{B}) \leq \Nprim$.

\begin{theorem}[Random Oracle Multi-User Indistinguishability]
\label{thm:mu-ind-ro}
For every flat-RO mu-ind adversary $\mathcal{B}$ with
execution-uniform user cap $D$ and resource counts of
\Cref{sec:mu-setting},
\[
  \Adv^{\mathsf{mu\text{-}ind}}_{\RO}(\mathcal{B})
  \;\leq\;
  \frac{D(D - 1)}{2^{k + 1}}
  + D \cdot \KeyGuess_k(\qRO(\mathcal{B}))
  + \frac{\nleaf(\mathcal{B})\bigl(\nleaf(\mathcal{B}) - 1\bigr)}
         {2^{\tauleaf + 1}}
  + \frac{\qv(\mathcal{B})}{2^{\tauroot}}.
\]
Here $\qv$ is the non-replay verification-query count of
\Cref{sec:resources}.
\end{theorem}

\begin{proof}
Use the pre-sampled-key formulation of \Cref{sec:mu-setting}: sample
keys $K_1,\dots,K_D$ at the start of the experiment, and let the $d$-th
$\mathsf{New}$ call activate $K_d$, unused keys being ignored.
Let $\mathsf{coll}$ denote the event that two of the pre-sampled keys
$K_1,\dots,K_D$ collide;
$\Pr[\mathsf{coll}] \leq \binom{D}{2}/2^k = D(D-1)/2^{k+1}$ by union
bound. In the random oracle model, use a standard hybrid over the $D$
users: $H_d$ answers users $1,\dots,d$ with $(\Rand,\Replay)$ and
users $d+1,\dots,D$ with the real construction interface, so $H_0$ is
the real mu-ind game and $H_D$ is the ideal one. Let $H_d^\star$ be
$H_d$ with a fixed output bit on $\mathsf{coll}$. Since
$\mathsf{coll}$ has the same probability in all hybrids,
\[
  \Adv^{\mathsf{mu\text{-}ind}}_{\RO}(\mathcal{B})
  \leq
  \Pr[\mathsf{coll}]
  +
  \bigl|\Pr[\mathcal{B}^{H_0^\star} \Rightarrow 1]
       - \Pr[\mathcal{B}^{H_D^\star} \Rightarrow 1]\bigr|.
\]
By the triangle inequality, it remains to bound each adjacent stopped
hybrid. For each $d$, \Cref{lem:mu-adjacent-hybrid} gives
\[
  \bigl|\Pr[\mathcal{B}^{H_{d-1}^\star} \Rightarrow 1]
       - \Pr[\mathcal{B}^{H_d^\star} \Rightarrow 1]\bigr|
  \leq
  \KeyGuess_k(\qRO(\mathcal{B}))
  + \frac{\nleaf^{(d)}(\nleaf^{(d)}-1)}{2^{\tauleaf+1}}
  + \frac{\qv^{(d)}}{2^{\tauroot}}.
\]

Summing across $d = 1,\dots,D$ gives
$D \KeyGuess_k(\qRO(\mathcal{B}))$ for the key-test contribution. The
leaf-collision contribution satisfies
\[
  \sum_d \frac{\nleaf^{(d)}(\nleaf^{(d)} - 1)}{2^{\tauleaf + 1}}
  = \frac{1}{2^{\tauleaf}} \sum_d \binom{\nleaf^{(d)}}{2}
  \leq \frac{1}{2^{\tauleaf}} \binom{\sum_d \nleaf^{(d)}}{2}
  = \frac{\nleaf(\mathcal{B})\bigl(\nleaf(\mathcal{B}) - 1\bigr)}
          {2^{\tauleaf + 1}},
\]
because the unordered pairs drawn within the per-user blocks are a
subset of all unordered pairs among $\sum_d \nleaf^{(d)} =
\nleaf(\mathcal{B})$ items. The verification-tag contribution satisfies
\[
  \sum_d \qv^{(d)} / 2^{\tauroot}
  = \qv(\mathcal{B}) / 2^{\tauroot}.
\]
Combining these estimates with the $\mathsf{coll}$ penalty gives the
stated RO bound.
\end{proof}

\begin{theorem}[Multi-User Indistinguishability]
\label{thm:mu-ind}
For every mu-ind adversary $\mathcal{A}$ against \TW{} with
execution-uniform user cap $D$ and the global resources of
\Cref{sec:mu-setting},
\[
  \Advmuind_\TW(\mathcal{A})
  \;\leq\;
  \Lift_c(\Ntot)
  + \frac{D(D - 1)}{2^{k + 1}}
  + D \cdot \KeyGuess_k(\Nprim)
  + \frac{\Nleaf(\Nleaf - 1)}{2^{\tauleaf + 1}}
  + \frac{\Qv}{2^{\tauroot}}.
\]
\end{theorem}

\begin{proof}
By the Lift Application corollary (\Cref{cor:lift-application}) for the
mu-ind game, there is a random oracle mu-ind adversary
$\mathcal{B} = \mathcal{B}_{\mathrm{derived}}(\mathcal{A})$ with
$\qRO(\mathcal{B}) \leq \Nprim$, $\nleaf(\mathcal{B}) \leq \Nleaf$,
$\qv(\mathcal{B}) \leq \Qv$, and the per-user construction-side caps of
\Cref{sec:mu-setting} such that
\[
  \Advmuind_\TW(\mathcal{A})
  \;\leq\;
  \Lift_c(\Ntot) + \Adv^{\mathsf{mu\text{-}ind}}_{\RO}(\mathcal{B}).
\]
\Cref{thm:mu-ind-ro} bounds the random oracle term. Substituting
$\qRO(\mathcal{B}) \leq \Nprim$,
$\nleaf(\mathcal{B}) \leq \Nleaf$, and
$\qv(\mathcal{B}) \leq \Qv$ gives the stated bound.
\end{proof}

\begin{corollary}[All-in-One AE Security]
\label{cor:ae}
For every nonce-respecting AE adversary $\mathcal{A}$ against \TW{}
with the resources of \Cref{sec:resources},
\[
  \Advae_\TW(\mathcal{A})
  \;\leq\;
  \Lift_c(\Ntot)
  + \KeyGuess_k(\Nprim)
  + \frac{\Nleaf(\Nleaf - 1)}{2^{\tauleaf + 1}}
  + \frac{\Qv}{2^{\tauroot}}.
\]
\end{corollary}

\begin{proof}
Build a mu-ind adversary $\mathcal{A}^{\mu}$ that makes one
$\mathsf{New}$ query, runs $\mathcal{A}$, and forwards each
encryption and verification query of $\mathcal{A}$ to user $1$. The
primitive interface is forwarded unchanged. The real and ideal views of
$\mathcal{A}$ in the all-in-one AE experiment are exactly the real and
ideal views of $\mathcal{A}^{\mu}$ in the mu-ind experiment with
$D = 1$, and the resource counts are unchanged. Applying
\Cref{thm:mu-ind} with $D = 1$ eliminates the key-collision term and
gives the stated bound.
\end{proof}

\begin{theorem}[Random Oracle \textsf{INT-RUP} Integrity]
\label{thm:int-rup-ro}
For every flat-RO \textsf{INT-RUP} adversary $\mathcal{B}$ against \TW{}
with nonce-respecting encryption queries and the resources of
\Cref{sec:resources},
\[
  \AdvintRUP_{\RO}(\mathcal{B})
  \;\leq\;
  \KeyGuess_k(\qRO(\mathcal{B}))
  + \frac{\nleaf(\mathcal{B})\bigl(\nleaf(\mathcal{B}) - 1\bigr)}
         {2^{\tauleaf + 1}}
  + \frac{\qv(\mathcal{B})}{2^{\tauroot}}.
\]
Plaintext-computing queries contribute to $N$ and $\nleaf$, but not to
$\qv$.
\end{theorem}

\begin{proof}
Work in the flat-RO \textsf{INT-RUP} game of \Cref{sec:ro-model}. Let
$\forge$ be the event that a valid non-replay verification query accepts.
Let $\badkey$ be the event that a direct $\ROflat$ query decodes under
$\KeyedTWOut$ to the hidden key $K$ and a counted output-point class, and let
$\badleaf$ be the event that two distinct construction-generated final leaf
transcripts receive the same $\tauleaf$-bit leaf tag projection. Final leaf
transcripts generated by encryption, plaintext-computing, and verification
computations are all included in $\nleaf$.

The key-test term is unchanged by the plaintext-computing oracle. The
key-renaming argument of \Cref{lem:ver-key-uniformity} applies with plaintext
oracle replies added to the visible prefix: under the online renaming
$K_0 \mapsto K_1$, every construction lookup that determines a returned
ciphertext, plaintext block, leaf tag, or root tag is paired with a distinct
renamed lookup carrying the same sampled value. A plaintext reply is obtained
by XORing the submitted ciphertext block with the same coupled keystream value,
and the recomputed tag is not returned. Thus, before $\badkey$, the hidden key
remains uniform outside the candidate keys already tested by direct queries,
and \Cref{lem:key-tests} gives
\[
  \Pr[\badkey] \leq \KeyGuess_k(\qRO(\mathcal{B})).
\]

For the leaf term, before $\badkey$ no direct query has pre-read a hidden-key
final leaf tag point: such a query would decode under $\KeyedTWOut$ to $K$ and
the corresponding $\Ltag(j)$ class. The hidden-key form of
\Cref{lem:leaf-collision}, applied to the enlarged set of construction
computations, gives
\[
  \Pr[\badleaf \text{ before } \badkey]
  \leq
  \frac{\nleaf(\mathcal{B})\bigl(\nleaf(\mathcal{B}) - 1\bigr)}
       {2^{\tauleaf + 1}}.
\]

It remains to bound accepting verification before $\badkey \cup \badleaf$.
Consider the stopped execution in which neither event has occurred. A final
root transcript class is \emph{tag-revealed} once an encryption query has
returned the root tag for that class. If a non-replay verification submission
belongs to a tag-revealed class, then by final-root determinacy
(\Cref{lem:transcript-inj}) and Leaf-Transcript Recovery
(\Cref{lem:leaf-recovery}) it has the same $(U,A,C)$ as the encryption query
that revealed the class. Submitting the revealed tag is a replay-table hit, and
submitting any other tag fails the comparison. Hence a successful non-replay
verification must belong to a class whose root tag has not yet been revealed by
encryption.

For such unrevealed classes, use the all-reject continuation of
\Cref{lem:ver-first-root-guess}, augmented with plaintext-computing queries.
Encryption queries and plaintext-computing queries are answered normally, and
each non-replay verification query performs the same construction computation
but returns reject in the continuation. Plaintext-computing queries may generate
the same stream and leaf tag transcripts as decryption, but they do not reveal
the final root tag value. Equivalently, the final root tag sample for an
unrevealed class can be deferred until it is first needed for an encryption
answer or a verification comparison; this does not change any plaintext reply.
By Output-Point Separation (\Cref{cor:output-sep}), the returned plaintext
blocks are projections at stream output points, not at the final root tag output
point.

Fix an unrevealed final root transcript class $\mathcal{C}$ and let
$G_{\mathcal{C}}$ be the set of distinct tags submitted by non-replay
verification queries in that class before the first later encryption query that
reveals the class, if any. Before that cutoff, no encryption query has returned
the class tag and no direct query has read the hidden-key root tag point without
triggering $\badkey$. Conditioned on the all-reject path and on all lazy-table
values except this final root tag projection, the correct tag for
$\mathcal{C}$ is uniform in $\bits^{\tauroot}$ and independent of
$G_{\mathcal{C}}$. The probability that any pre-cutoff verification submission
in the class accepts is therefore at most
$|G_{\mathcal{C}}|/2^{\tauroot}$.

Each valid non-replay verification query contributes its submitted tag to at
most one such set $G_{\mathcal{C}}$, and replay-table hits do not contribute.
Thus $\sum_{\mathcal{C}} |G_{\mathcal{C}}| \leq \qv(\mathcal{B})$, and a union
bound over classes gives
\[
  \Pr[\forge \text{ before } \badkey \cup \badleaf]
  \leq \qv(\mathcal{B}) / 2^{\tauroot}.
\]
Combining the three bounds proves the theorem.
\end{proof}

\begin{theorem}[\textsf{INT-RUP} Integrity]
\label{thm:int-rup}
For every \textsf{INT-RUP} adversary $\mathcal{A}$ against \TW{} with
nonce-respecting encryption queries and the resources of \Cref{sec:resources},
\[
  \AdvintRUP_\TW(\mathcal{A})
  \;\leq\;
  \Lift_c(\Ntot)
  + \KeyGuess_k(\Nprim)
  + \frac{\Nleaf(\Nleaf - 1)}{2^{\tauleaf + 1}}
  + \frac{\Qv}{2^{\tauroot}}.
\]
\end{theorem}

\begin{proof}
By the Lift Application corollary (\Cref{cor:lift-application}) for the
\textsf{INT-RUP} game, there is a random oracle \textsf{INT-RUP} adversary
$\mathcal{B} = \mathcal{B}_{\mathrm{derived}}(\mathcal{A})$ with
$\qRO(\mathcal{B}) \leq \Nprim$, $\nleaf(\mathcal{B}) \leq \Nleaf$, and
$\qv(\mathcal{B}) \leq \Qv$ such that
\[
  \AdvintRUP_\TW(\mathcal{A})
  \;\leq\;
  \Lift_c(\Ntot) + \AdvintRUP_{\RO}(\mathcal{B}).
\]
\Cref{thm:int-rup-ro} bounds the random oracle term, and substituting the
derived-adversary resource caps gives the stated bound.
\end{proof}
